\section{Results}
\label{sec:results}

\subsection{Forward accuracy}
\label{sec:forward}

Table~\ref{tab:metrics} compares the model with the five classical correlations
on the held-out test split. The PINN attains a $\dTonb$ RMSE of \SI{1.77}{K}
($R^2=0.79$) and a $\qonb$ RMSE of \SI{133}{kW\,m^{-2}} ($R^2=0.78$). The best
classical correlation (Bergles--Rohsenow) reaches \SI{3.73}{K} on $\dTonb$, and
all correlations have strongly negative $R^2$ on $\qonb$---i.e.\ worse than
predicting the mean---because they were not built for this combined parameter
range. The PINN is thus about $2.1$--$2.7\times$ more accurate on $\dTonb$ and
more than an order of magnitude more accurate on $\qonb$. Figure~\ref{fig:parity}
shows the test parity on logarithmic axes: predictions lie within the
$\pm20\%$ band across the full $2$--\SI{30}{K} and $30$--\SI{1500}{kW\,m^{-2}}
ranges, including the low-superheat microchannel and high-pressure points.
Across sources (Fig.~\ref{fig:persource}) the residual error is concentrated
rather than uniform. The large modern sweeps that anchor the corpus are fit
well---Cheng ($n=54$, \SI{0.60}{K}) and Kuang ($n=66$, \SI{0.89}{K}), with Liu
($n=23$, \SI{1.78}{K}) at the overall level---as are the smaller well-instrumented
water sets (Wang, Hong, Al-Yahia; $0.5$--\SI{0.8}{K}). The larger per-source
errors fall on the sparse literature sets in which only a handful of records
report wall superheat: Qu ($n=11$, \SI{3.0}{K}), Forrest ($n=14$, \SI{4.7}{K})
and the two- and three-point Bergles and Kandlikar sets ($4.6$--\SI{5.9}{K}),
where a few digitised incipience points dominate the RMSE. The densely sampled
operating envelope is therefore captured accurately, with the residual error
confined to low-count, heat-flux-dominated subsets.

\begin{figure}[t]
\centering
\includegraphics[width=0.72\linewidth]{fig_per_source_rmse.png}
\caption{Per-source ONB wall-superheat RMSE (PINN, all data); $n$ is the number
of points with a measured $\dTonb$ (a subset of each source's total rows in
Table~\ref{tab:sources}). Dashed line: the overall held-out test RMSE of
\SI{1.77}{K}.}
\label{fig:persource}
\end{figure}

\begin{table}[t]
\centering
\caption{Test-split accuracy: PINN versus classical correlations (same
$n_{\Delta T}=36$, $n_q$ varies by correlation applicability).}
\label{tab:metrics}
\small
\begin{tabular}{lcccc}
\toprule
Method & RMSE$_{\dTonb}$ (K) & $R^2_{\dTonb}$ & RMSE$_{\qonb}$ (kW\,m$^{-2}$) & $R^2_{\qonb}$ \\
\midrule
\textbf{PINN (this work)}   & \textbf{1.77} & \textbf{0.79} & \textbf{133}  & \textbf{0.78} \\
Bergles--Rohsenow \cite{BERGLES1964} & 3.73 & 0.01 & 1670 & $-59$ \\
Hsu \cite{HSU1962}                   & 4.49 & $-0.44$ & 2745 & $-162$ \\
Sato--Matsumura \cite{SATOMATSUMURA1964}      & 4.49 & $-0.44$ & 2745 & $-162$ \\
Kandlikar \cite{KANDLIKAR1991}       & 4.49 & $-0.44$ & 2745 & $-162$ \\
Basu et al.\ \cite{BASU2002}         & 4.78 & $-0.63$ & 3366 & $-244$ \\
\bottomrule
\end{tabular}
\end{table}

\begin{figure}[t]
\centering
\includegraphics[width=0.49\linewidth]{parity_dT_onb_test.png}
\includegraphics[width=0.49\linewidth]{parity_q_onb_test.png}
\caption{Test-split parity on logarithmic axes for (a)~the ONB wall superheat
$\dTonb$ and (b)~the ONB heat flux $\qonb$. Dashed line: ideal ($1{:}1$); shaded
band: $\pm20\%$.}
\label{fig:parity}
\end{figure}

\subsection{The role of transfer learning and data curation}

Accuracy improved monotonically as the dataset was cleaned and extended and as
two physical fixes were applied (Section~\ref{sec:physcon}). Removing
pool-boiling contamination (Section~\ref{sec:data}) was the single largest
gain, cutting the $\dTonb$ RMSE from \SI{4.0}{K} to \SI{1.85}{K}; adding the
microchannel data sharpened $\qonb$ ($R^2$ from $0.54$ to $0.78$); and the
high-pressure anchors with the reduced-pressure feature lowered $\dTonb$ RMSE to
\SI{1.84}{K} while fixing the pressure trend; a final regime-dependent Nusselt
closure in the Hsu constraint (Section~\ref{sec:discussion}) then brought the
test error to \SI{1.77}{K} ($R^2=0.79$; stage S9 in Fig.~\ref{fig:ablation}).

A controlled ablation isolates the role of the transferred surface encoder
(Table~\ref{tab:ablation}). A surface encoder trained from random
initialisation matches the frozen transferred one on wall superheat
(\SI{1.77}{K} vs \SI{1.84}{K}), so transfer learning is not required for raw
accuracy. Its value is rather that of an \emph{informative frozen prior}: the
frozen transferred encoder outperforms a frozen \emph{random} encoder by about
\SI{0.5}{K} (\SI{1.84}{K} vs \SI{2.36}{K}) at no additional trainable cost. This
matters because most flow-boiling studies omit the surface descriptors $R_a$ and
$\theta$, so a frozen prior that already encodes a sensible cavity-activation
representation regularises the model where those inputs are missing.

\begin{figure}[t]
\centering
\includegraphics[width=0.82\linewidth]{fig_ablation_progression.png}
\caption{Test accuracy across the successive model stages S1--S8. Bars:
$\dTonb$ RMSE; line: $R^2$; annotations mark the change at each step.
Wall-superheat error falls 77\% from the transfer baseline (S1).}
\label{fig:ablation}
\end{figure}

\begin{table}[t]
\centering
\caption{Transfer-learning ablation (same test split, $n_{\dTonb}=36$), run
under a common Dittus--Boelter Hsu closure to isolate the encoder. A trained
encoder matches the frozen transferred one; the transferred prior beats a frozen
\emph{random} encoder. (The adopted model additionally uses the regime-dependent
Nu of Section~\ref{sec:discussion}, which shifts these values by $<\!0.1$~K.)}
\label{tab:ablation}
\small
\begin{tabular}{llcc}
\toprule
Variant & Surface encoder & RMSE$_{\dTonb}$ (K) & $R^2_{\dTonb}$ \\
\midrule
A (this work) & transferred, frozen   & \textbf{1.84} & 0.77 \\
B            & random init, trained  & 1.77 & 0.78 \\
C            & random init, frozen   & 2.36 & 0.62 \\
\bottomrule
\end{tabular}
\end{table}

\subsection{Cross-study generalization}
\label{sec:loso}

The headline test error above uses a per-study stratified split, in which every
study contributes to both training and test; it therefore measures interpolation
\emph{within} the sampled studies. To probe genuine
study-to-study extrapolation we retrain under leave-one-study-out (LOSO): one
study is held out \emph{entirely} as the test set while the model trains on the
other nine (Table~\ref{tab:loso}). Accuracy degrades sharply---to
$5.7$--\SI{10.9}{K} with strongly negative $R^2$---for every held-out study,
including a central, low-pressure one (Cheng). The model thus interpolates
reliably within the trained envelope but does not extrapolate to an entirely
unseen study. We regard interpolation within the envelope as its intended
operating regime and return to this limitation in
Section~\ref{sec:discussion}.

\begin{table}[t]
\centering
\caption{Leave-one-study-out generalization: each study is held out entirely as
the test set and the model trains on the remaining nine (Dittus--Boelter Hsu
closure, matching the holdout runs). The within-envelope stratified split is
shown for reference.}
\label{tab:loso}
\small
\begin{tabular}{llcc}
\toprule
Held-out study & Regime & RMSE$_{\dTonb}$ (K) & $R^2_{\dTonb}$ \\
\midrule
Wang \cite{WANG2024}          & high pressure, 10--16\,MPa & 5.70  & $-87$ \\
Qu \& Mudawar \cite{QUMUDAWAR2002} & microchannel, 0.35\,mm     & 7.70  & $-1.5$ \\
Liu \cite{LIU2005}            & microchannel               & 6.87  & $-2.9$ \\
Cheng \cite{CHENG2022}        & narrow channel             & 10.9  & $-227$ \\
\midrule
\textbf{Stratified split (ref.)} & all studies            & \textbf{1.84} & \textbf{0.77} \\
\bottomrule
\end{tabular}
\end{table}

\subsection{Physics consistency}
\label{sec:physcon}

A predictor intended for design must respect the established physical trends.
We verify the sign of $\partial\dTonb/\partial\xi$ for each governing variable.

\subsubsection{Subcooling--Hsu conflict (low superheat)}
\label{sec:patternB}
An early model systematically under-predicted $\dTonb$ in the low-superheat
($<\!\SI{5}{K}$), high-subcooling regime. The cause was the Hsu
\emph{equality} coupling: at the low heat fluxes of these points the
saturated-pool Hsu superheat is near zero ($\sim\!\SI{0.06}{K}$), and the
equality term dragged predictions down against measured values of
$2$--\SI{5}{K}. Replacing it with the one-sided lower bound of
Eq.~\eqref{eq:onesided} removed the bias: on the most affected source the test
error became unbiased (mean error $+\SI{0.03}{K}$, RMSE \SI{0.62}{K}). This
confirms that the saturated-pool criterion should bound, not equal, the
subcooled-flow superheat.

\subsubsection{Pressure dependence}
\label{sec:patternC}
Pressure has a first-order effect ($P\!\uparrow\!\Rightarrow\!\dTonb\!\downarrow$,
the Bergles--Rohsenow trend), yet a model in which pressure entered only through
fluid properties was \emph{blind} to it: the correlation between pressure and
predicted superheat was $+0.01$ against a measured $-0.14$ over the
high-pressure subset. Adding reduced pressure as an explicit input feature
flipped the sign of the predicted correlation, and adding the
\SI{10}{}--\SI{16}{MPa} anchors strengthened it to $-0.69$ (measured $-0.57$
over the combined high-pressure data); the model then tracks the Wang et al.\
pressure sweep in the correct descending direction (Fig.~\ref{fig:ptrend}).
Property-based nondimensionalization alone is therefore insufficient to encode
the pressure dependence of incipience; an explicit reduced-pressure input is
required.

\begin{figure}[t]
\centering
\includegraphics[width=0.72\linewidth]{fig_pressure_trend.png}
\caption{Measured (solid) and predicted (dashed) ONB wall superheat versus
pressure for two fixed-condition sweeps (Kuang 2024: $G=500$, $q''=300$;
Wang 2024: $G=1000$, $q''=200$; $G$ in kg\,m$^{-2}$s$^{-1}$, $q''$ in
kW\,m$^{-2}$). A residual over-prediction persists in the $3.5$--\SI{5}{MPa}
single-source pocket.}
\label{fig:ptrend}
\end{figure}

\subsubsection{Remaining trends}
On representative sweeps the model satisfies $\dTsub\!\uparrow\!\Rightarrow\!
\dTonb\!\uparrow$ (correct sign, modest magnitude) and, through the Hsu coupling
and the data, $q\!\uparrow\!\Rightarrow\!\dTonb\!\uparrow$ (in-source
correlations $+0.57$ to $+0.88$); both are shown in Fig.~\ref{fig:trends}. The
subcooling sensitivity is weak (a fraction of a kelvin over the range), whereas
the mass-flux---and hence Reynolds-number---dependence is itself
pressure-dependent: modest at elevated pressure but appreciable at low pressure
(a $\dTonb$ swing of several kelvin across $G=200$--\SI{2000}{kg\,m^{-2}s^{-1}}),
where a higher mass flux raises the wall superheat required for incipience---the
expected flow-convection trend, visible in the design map of
Fig.~\ref{fig:designmap}.

\begin{figure}[t]
\centering
\includegraphics[width=\linewidth]{fig_physics_trends.png}
\caption{Physical-trend consistency. (a)~Model prediction increases
monotonically with subcooling---correct sign but small magnitude (a fraction of
a kelvin, as for Reynolds number and mass flux). (b)~Measured ONB superheat
rises with heat flux across sources.}
\label{fig:trends}
\end{figure}

\subsection{Uncertainty quantification}
\label{sec:uq}

A five-member deep ensemble (identical data split, varied initialisation)
provides predictive uncertainty. The ensemble mean is also \emph{more accurate}
than any single model ($\dTonb$ RMSE \SI{1.72}{K}, $R^2=0.80$;
Table~\ref{tab:metrics} single-model value \SI{1.77}{K}), a known benefit of
ensembling. Figure~\ref{fig:uq} shows the ensemble parity with $\pm2\sigma$
epistemic bars, the calibration curve, and the epistemic spread versus pressure.

The calibration analysis is deliberately honest. The mean epistemic standard
deviation is \SI{0.51}{K}, and the $1\sigma$/$2\sigma$ coverages are only
$28\%$/$42\%$ (ideal $68\%$/$95\%$): the ensemble \emph{under}-estimates total
predictive uncertainty by roughly $4.7\times$. This is expected---deep-ensemble
spread captures only epistemic (model) disagreement, whereas the dominant
uncertainty here is aleatoric: digitisation error and genuine experimental
scatter of order \SI{1}{K}. A single scaling factor of $4.7$ yields
well-calibrated intervals as an interim measure (green curve,
Fig.~\ref{fig:uq}b). The principled alternative---a heteroscedastic variance
head trained by negative log-likelihood---was implemented and tested, but at the
present data scale it did not improve calibration ($1\sigma$ coverage unchanged
at $28\%$) and slightly degraded accuracy, indicating that a reliable per-point
aleatoric variance is not learnable from $\sim\!240$ heterogeneous points
(Section~\ref{sec:discussion}). Notably the epistemic spread does \emph{not} flag the
high-pressure region as uncertain (the members agree there, anchored by the
Wang et al.\ data), confirming that the residual error there is
aleatoric/structural rather than a lack of model consensus.

\begin{figure}[t]
\centering
\includegraphics[width=\linewidth]{ensemble_uq.png}
\caption{Deep-ensemble ($K=5$) uncertainty. (a)~Test parity with $\pm2\sigma$
epistemic bars. (b)~Calibration: observed coverage versus interval half-width
$k$ ($\times\sigma$); the epistemic-only spread (pink) under-covers, a
$4.7\times$ rescaling (green) nears the Gaussian ideal (dashed). (c)~Epistemic
spread versus pressure.}
\label{fig:uq}
\end{figure}
